\section{Hodógrafo de Apophis}

\subsection{Hodógrafo kepleriano vs. perturbado}

Calcular el hodógrafo teórico ---la circunferencia con centro $(0, \mu e / h)$ y radio $\mu / h$--- y superponerlo sobre el hodógrafo numérico obtenido de la simulación. Mostrar que en 2 cuerpos coinciden perfectamente y que en $N$ cuerpos la circunferencia se deforma ligeramente.

\subsection{Hodógrafo durante el encuentro}

Hacer un acercamiento temporal al hodógrafo justo durante los días del encuentro cercano. El vector velocidad debería trazar una curva distinta de la circunferencia kepleriana usual, precisamente cuando la Tierra modifica de forma apreciable la dinámica.

\subsection{Verificación analítica del hodógrafo}

Dado el vector de estado inicial de Apophis, calcular analíticamente el centro y el radio del hodógrafo, $(0, \mu e / h)$ y $\mu / h$, y verificar numéricamente que todos los puntos $(v_x, v_y)$ de la simulación caen sobre esa circunferencia en el caso de 2 cuerpos.